\section{A personal note}

What follows is a more personal view of bike-sharing and of the application of \ac{ai} in these systems. I am not a regular user, as I have only used \acp{bss} a few times, but I have watched the impact of \textit{Bicing} on Barcelona since it was launched in 2007.

I believe bike-sharing was, is, and will be a genuine step toward cleaner urban mobility. It began as a social and ecological initiative~\cite{vanderZee2016}, and the boom of the 2000s led many cities to adopt it. I think many of those cities did so mainly because it had become fashionable: a green and photogenic flagship, copied once another city had already shown that it could be done~\cite{Medard2019}. Even so, it is thanks to that trend that so many cities now have a \ac{bss}, and as these systems mature, they prove they can still fulfill that original promise.

In Barcelona, launching \textit{Bicing} in 2007 was one of the first widely visible steps toward a more cyclable city. It made cycling look ordinary in a place long organized around the car, and it gave people who would never have bought a bicycle a reason to try one. The city later continued in that direction with superblocks, low-emission zones, and a much more extensive cycling network. Thanks to those efforts together, I think a cycling culture is slowly taking hold, and that Barcelona is becoming a more pleasant place to live because of these choices.

As \acp{bss} matured in Barcelona and in other cities, research interest moved on to other topics. Further development of the systems is now led mostly by the companies and municipalities that run them, and academic work has become much scarcer than it was in the 2010s. There is still room to improve them, as there is in any transport network, but day-to-day operation leaves operators with little time or capacity to look beyond the next service day. That is why I think research can still help.

The same maturity that pulled research away is also what now makes that help possible. These systems record how bicycles are used, how they move, and how they fail---traces that barely existed when bike-sharing was still new---and the \ac{ai} models that can be built from them have improved a great deal in the last few years. Those traces could be used in many ways; to give only two examples, they might be joined with app searches to show demand that never appears as a trip, or with mobile-phone or \ac{pt} card data to show which mode a ride replaced. The same logic applies, perhaps more strongly, to dockless systems, which this thesis does not examine: they record continuous \ac{gps} traces of the routes themselves, not only origins and destinations, and that would allow a closer look at how people actually ride.

Even so, the most complete data and the most capable models would not, by themselves, improve these systems. What matters most is a political and operational will to keep making the effort: a directed and continuous desire to improve the service so that it reaches as many people as it can. That means a service that is affordable, and one that works---bicycles in good condition, docks with bikes and free spaces, batteries that last. I believe that will is present in Barcelona. The city has gone on expanding the network and replacing \acp{m-b} with \acp{e-b}, and those choices have made the system easier and more convenient to use. I am optimistic that it will last, and not only here. As more people care about a warming climate and about cities that are pleasant to live in, I think the will to improve these systems will remain. Research can still show where the service falls short, but acting on that is what closes the gap.
